\documentclass{article}
\usepackage{multicol}

\begin{document}
\begin{multicols}{2}


    
Notes

RM)M3.1.1 Managing risk

Learning goals:

1)  What is risk ?

2)  How can we calculate risk ?

3)  How can we mitigate the risk ?

Introduction

**Risky** versus **Uncertain events**

Risky:

Likelihood / probability of the outcomes can be [calculated]{.underline}.\
Through past trials or other methods

Ex: A coin toss (H or T)

Uncertain:

[Cannot quantify]{.underline} the probability of the outcome.

Ex: Technological breakthroughs, Pandemics, etc

Statistical concepts

**Probability density function:**

It is a function that assigns the probability of the outcomes of risky events.

Can be discrete or continuous.

**Random variables:**

Numerical value to the outcome of a risky event ?

**Expected values:**

Mean of the outcomes.

Represented with$$
E(y)=\sum_{j=1}^n p_j y_j 
$$

Where

\space p~j -~ Probability of event j

\space Yj - Random variable j

AND

$$
\sum_{j=1}^n p_j = 1$$

is always true.

**Variance and standard deviation:**

Basic formula


Variance is used to calculate the "risk" fo the given random variable.

**Risk premiums and certainty equivalent incomes:**

Concept of **certainty equivalent income** - Is the amount of money you would take in [absolute certainty]{.underline} [instead]{.underline} of taking a [riskier trade]{.underline} with a potential for higher return.

Concept of **risk premium** -

If an investor would pay a [positive risk premium]{.underline} to receive an investments expected value with certainty, then the investor is [risk averse]{.underline}.

If the investor would [pay nothing]{.underline} to receive an investments expected value with certainty, the investor is [risk neutral]{.underline}.

If the investor [would pay]{.underline} to keep the variability (the investor enjoys gambling), we would say the investor is [risk preferring]{.underline}.

We can calculate the Certainty Equivalent Income:

CEI = E(x) - Premium

$$
y_{C E}^i=E(y)-\frac{\lambda^i}{2} \sigma_y^2
$$

y~ce~ - Certainty Equivalent Income

E(y) - Expected value of the investment

The other - Risk Premium

lambda - Risk preference measure

lambda/2 - The risk aversion coefficient (Imagine it like a slope)

**Expected value-variance criterion:**

\end{multicols}{2}  
\end{document}
